\documentclass[11pt]{article}
\usepackage[margin=0.9in]{geometry}
\usepackage{amsmath,amssymb,amsthm,mathtools}
\usepackage[T1]{fontenc}
\usepackage[utf8]{inputenc}
\usepackage{enumitem}
\usepackage{booktabs,longtable,array,seqsplit}
\usepackage{hyperref}
\usepackage{xcolor}
\usepackage{microtype}
\hypersetup{colorlinks=true,linkcolor=blue!45!black,citecolor=blue!45!black,urlcolor=blue!45!black}
\setlength{\emergencystretch}{6em}

\newtheorem{theorem}{Theorem}[section]
\newtheorem{proposition}[theorem]{Proposition}
\newtheorem{lemma}[theorem]{Lemma}
\newtheorem{corollary}[theorem]{Corollary}
\newtheorem{claim}[theorem]{Claim}
\theoremstyle{definition}
\newtheorem{definition}[theorem]{Definition}
\newtheorem{assumption}[theorem]{Assumption}
\theoremstyle{remark}
\newtheorem{remark}[theorem]{Remark}

\newcommand{\R}{\mathbb{R}}
\newcommand{\C}{\mathbb{C}}
\newcommand{\N}{\mathbb{N}}
\newcommand{\Hsp}{\mathcal{H}}
\newcommand{\Fsp}{\mathcal{F}}
\newcommand{\Id}{\mathrm{Id}}
\newcommand{\tr}{\operatorname{tr}}
\newcommand{\diag}{\operatorname{diag}}
\newcommand{\rank}{\operatorname{rank}}
\newcommand{\spec}{\operatorname{spec}}
\newcommand{\spanop}{\operatorname{span}}
\newcommand{\dist}{\operatorname{dist}}
\newcommand{\range}{\operatorname{ran}}
\newcommand{\Ker}{\operatorname{ker}}
\newcommand{\Pinch}{\mathcal{C}}
\newcommand{\Off}{\widetilde{\mathcal{C}}}
\newcommand{\Pperp}{P^{\perp}}
\newcommand{\Qperp}{Q^{\perp}}
\newcommand{\op}{\mathrm{op}}
\newcommand{\F}{\mathrm{F}}
\newcommand{\HS}{\mathrm{HS}}
\newcommand{\ip}[2]{\left\langle #1,#2\right\rangle}
\newcommand{\norm}[1]{\left\lVert #1\right\rVert}
\newcommand{\abs}[1]{\left\lvert #1\right\rvert}
\newcommand{\code}[1]{\nolinkurl{#1}}

\newenvironment{argument}{\paragraph{Argument.}\begin{enumerate}[leftmargin=2em]}{\end{enumerate}}
\newenvironment{proofoutline}{\paragraph{Proof architecture.}\begin{enumerate}[leftmargin=2em]}{\end{enumerate}}
\newcommand{\sourceanchor}[1]{\par\smallskip\noindent\textbf{Source anchor.} #1\par\smallskip}
\newcommand{\sourcecomment}[1]{\begin{remark}#1\end{remark}}
\newenvironment{sourceguide}{\begin{description}[style=nextline,leftmargin=3.3cm,labelwidth=3.1cm]}{\end{description}}
\newenvironment{formalizationmap}{\begin{longtable}{@{}p{0.20\textwidth}p{0.31\textwidth}p{0.40\textwidth}@{}}\toprule Source item & Modern mathematical content & Repository status \\\midrule\endfirsthead\toprule Source item & Modern mathematical content & Repository status \\\midrule\endhead}{\bottomrule\end{longtable}}

\title{Source-order modernization of Yu--Wang--Samworth: a population-gap Davis--Kahan theorem for statistics}
\author{}
\date{2015 journal paper; reconstructed from the supplied 2014 preprint transcription}
\begin{document}
\maketitle

\section*{Source map and scope}
\begin{sourceguide}
\item[Primary source] Yi Yu, Tengyao Wang, and Richard J. Samworth, \emph{A Useful Variant of the Davis--Kahan Theorem for Statisticians}, Biometrika 102 (2015), 315--323; preprint arXiv:1405.0680.
\item[Local evidence] Supplied transcription {\small\code{Yu-Wang-Samworth-2014-Useful-DavisKahanVariant(1).tex}}.\\ SHA--256 \texttt{\seqsplit{ed45f19eb1a619a182a737aa971b2120727d2bebb86760852f860a9fbe54ecfe}}.
\item[Redistribution] The transcription is not copied into the repository. This note is a new, source-order mathematical reconstruction.
\item[Coverage] Principal-angle definitions; the statistical Davis--Kahan baseline; Theorem 2; sharpness examples; the rank-one corollary; the application audit; Theorem 3 for singular subspaces; Appendix Lemma 1; all displayed proof identities and constants.
\item[Formalization target] \code{YuWangSamworth.lean}, \code{AlignedBasis.lean}, \code{DavisKahanTheory/Statistics.lean}, and \code{DavisKahanTheory/SingularSubspace.lean}.
\end{sourceguide}

The purpose is not to repeat the paper's prose. The reconstruction isolates every mathematical assertion, normalizes notation, expands all hidden matrix identities, and records exactly how the paper's finite-dimensional proof maps to the intrinsic Lean development.

\section{Principal angles and distances between subspaces}

Let $V,\widehat V\in\R^{p\times d}$ have orthonormal columns:
\[
  V^TV=I_d,
  \qquad
  \widehat V^T\widehat V=I_d.
\]
Let
\[
  \sigma_1\ge\cdots\ge\sigma_d\ge0
\]
be the singular values of $\widehat V^TV$. The principal angles are
\[
  \theta_j:=\arccos\sigma_j\in[0,\pi/2].
\]
Write
\[
  \Theta(\widehat V,V):=\diag(\theta_1,\dots,\theta_d),
  \qquad
  \sin\Theta(\widehat V,V):=\diag(\sin\theta_1,\dots,\sin\theta_d).
\]
The Frobenius sine distance is
\[
  \norm{\sin\Theta(\widehat V,V)}_{\F}^2
  =\sum_{j=1}^d\sin^2\theta_j
  =d-\norm{\widehat V^TV}_{\F}^2.
\]
If $V_1$ is an orthonormal complement to $V$, then
\[
  \norm{V_1^T\widehat V}_{\F}^2
  =\norm{\sin\Theta(\widehat V,V)}_{\F}^2.
\]

The basis-sensitive distance is minimized over orthogonal alignments:
\[
  \min_{O\in O(d)}\norm{\widehat VO-V}_{\F}.
\]
If $O$ is the Procrustes alignment obtained from an SVD of $\widehat V^TV$, then
\[
  \norm{\widehat VO-V}_{\F}^2
  =2\sum_{j=1}^d(1-\cos\theta_j)
  \le2\sum_{j=1}^d\sin^2\theta_j.
\]
Therefore
\[
  \norm{\widehat VO-V}_{\F}
  \le\sqrt2\norm{\sin\Theta(\widehat V,V)}_{\F}.
\]

\section{The mixed-gap Davis--Kahan baseline used by statisticians}

\sourceanchor{Introduction, stated Davis--Kahan theorem and equation (1).}

Let $\Sigma,\widehat\Sigma\in\R^{p\times p}$ be symmetric, with decreasing eigenvalues
\[
  \lambda_1\ge\cdots\ge\lambda_p,
  \qquad
  \widehat\lambda_1\ge\cdots\ge\widehat\lambda_p.
\]
Fix $1\le r\le s\le p$, set $d=s-r+1$, and choose orthonormal eigenvector blocks
\[
  V=(v_r,\dots,v_s),
  \qquad
  \widehat V=(\widehat v_r,\dots,\widehat v_s),
\]
with
\[
  \Sigma v_j=\lambda_jv_j,
  \qquad
  \widehat\Sigma\widehat v_j=\widehat\lambda_j\widehat v_j.
\]
The source defines a mixed population--sample separation
\[
  \delta
  :=\inf\left\{
      \abs{\widehat\lambda-\lambda}:
      \lambda\in[\lambda_s,\lambda_r],\ 
      \widehat\lambda\in
      (-\infty,\widehat\lambda_{s-1}]
      \cup[\widehat\lambda_{r+1},\infty)
    \right\},
\]
with endpoint conventions printed as
\[
  \widehat\lambda_0=-\infty,
  \qquad
  \widehat\lambda_{p+1}=\infty.
\]

\sourcecomment{The interval indices and endpoint signs in the supplied transcription are awkward relative to decreasing eigenvalue order. This note records the source statement but does not treat those conventions as a new canonical formulation. The formal development uses explicit spectral-block separation predicates.}

\begin{theorem}[Statistical Davis--Kahan baseline]
If $\delta>0$, then
\[
  \norm{\sin\Theta(\widehat V,V)}_{\F}
  \le
  \frac{\norm{\widehat\Sigma-\Sigma}_{\F}}{\delta}.
\]
\end{theorem}
The source states that both Frobenius norms may be replaced simultaneously by the operator norm or, more generally, by any orthogonally invariant norm.

For a single eigenvector, $r=s=j$, this specializes to
\[
  \sin\Theta(\widehat v_j,v_j)
  \le
  \frac{\norm{\widehat\Sigma-\Sigma}_{\op}}
       {\min\bigl(
          \abs{\widehat\lambda_{j-1}-\lambda_j},
          \abs{\widehat\lambda_{j+1}-\lambda_j}
        \bigr)}.
\]
After choosing the sign of $\widehat v_j$ so that $\widehat v_j^Tv_j\ge0$,
\[
  \norm{\widehat v_j-v_j}
  \le\sqrt2\sin\Theta(\widehat v_j,v_j).
\]
The usual statistical workflow then invokes Weyl's inequality to replace the random mixed gaps by fractions of the population gaps on a high-probability event. The main theorem eliminates that intermediate event from the theorem statement.

\section{Main population-gap theorem}

\sourceanchor{Theorem 2, equations (2)--(3), sharpness discussion, equation (4), and Corollary 1.}

Define the population exterior gap around the block $r,\dots,s$ by
\[
  \Delta
  :=\min(\lambda_{r-1}-\lambda_r,\lambda_s-\lambda_{s+1}),
\]
with
\[
  \lambda_0:=+\infty,
  \qquad
  \lambda_{p+1}:=-\infty.
\]
Let
\[
  E:=\widehat\Sigma-\Sigma.
\]

\begin{theorem}[Yu--Wang--Samworth population-gap theorem]
If $\Delta>0$, then
\begin{equation}
  \norm{\sin\Theta(\widehat V,V)}_{\F}
  \le
  \frac{2\min\bigl(\sqrt d\norm{E}_{\op},\norm{E}_{\F}\bigr)}{\Delta}. \tag{YWS-1}
\end{equation}
Moreover, there exists $\widehat O\in O(d)$ such that
\begin{equation}
  \norm{\widehat V\widehat O-V}_{\F}
  \le
  \frac{2^{3/2}\min\bigl(\sqrt d\norm{E}_{\op},\norm{E}_{\F}\bigr)}{\Delta}. \tag{YWS-2}
\end{equation}
\end{theorem}

The paper emphasizes two improvements over the commonly quoted statistical Davis--Kahan form:
\begin{enumerate}[leftmargin=2em]
\item the denominator contains only population eigenvalues;
\item the numerator automatically chooses between the dimension-scaled operator norm and the Frobenius norm.
\end{enumerate}

Let
\[
  \Lambda:=\diag(\lambda_r,\dots,\lambda_s).
\]
The proof actually establishes the sharper residual forms
\[
  \norm{\sin\Theta(\widehat V,V)}_{\F}
  \le
  \frac{\norm{\widehat V\Lambda-\Sigma\widehat V}_{\F}}{\Delta},
\]
and
\[
  \norm{\widehat V\widehat O-V}_{\F}
  \le
  \frac{\sqrt2\norm{\widehat V\Lambda-\Sigma\widehat V}_{\F}}{\Delta}.
\]
The displayed theorem replaces the residual by a directly usable perturbation bound.

\subsection{Sharpness example with orthogonal target subspaces}

Let
\[
  \Sigma=\diag(\underbrace{3,\dots,3}_{d},
                       \underbrace{1,\dots,1}_{p-d}),
\]
and
\[
  \widehat\Sigma
  =\diag(\underbrace{2-\varepsilon,\dots,2-\varepsilon}_{p-d},
         \underbrace{2,\dots,2}_{d}),
\]
where $0<\varepsilon$ and $1\le d\le\lfloor p/2\rfloor$. The top-$d$ eigenspaces are orthogonal. Hence for every $O\in O(d)$,
\[
  \norm{\widehat VO-V}_{\F}
  =\sqrt2\norm{\sin\Theta(\widehat V,V)}_{\F}
  =\sqrt{2d}.
\]
The perturbation norm and population gap give
\[
  \sqrt{2d}
  \le\sqrt{2d}(1+\varepsilon)
  =
  \frac{2^{3/2}\sqrt d\norm{E}_{\op}}
       {\lambda_d-\lambda_{d+1}}.
\]
This shows that the aligned-basis constant and dimension dependence cannot be uniformly improved beyond the displayed constant scale.

\subsection{Sharpness scale for nearby one-dimensional eigenspaces}

Take
\[
  \Sigma=\diag(3,1),
\]
and
\[
  \widehat\Sigma
  =R_\varepsilon\diag(3,1)R_\varepsilon^T,
  \qquad
  R_\varepsilon=
  \begin{pmatrix}
    \sqrt{1-\varepsilon^2}&-\varepsilon\\
    \varepsilon&\sqrt{1-\varepsilon^2}
  \end{pmatrix}.
\]
For
\[
  v=(1,0)^T,
  \qquad
  \widehat v=(\sqrt{1-\varepsilon^2},-\varepsilon)^T,
\]
we have
\[
  \sin\Theta(\widehat v,v)=\varepsilon,
\]
\[
  \norm{\widehat v-v}^2
  =2-2\sqrt{1-\varepsilon^2},
\]
and
\[
  \frac{2\norm{\widehat\Sigma-\Sigma}_{\op}}{3-1}
  =2\varepsilon.
\]
Thus the sine bound is tight up to factor $2$, and the vector bound up to factor $2^{3/2}$, even in a small-angle regime.

\subsection{Why a double-angle estimate does not directly solve the statistical problem}

For oriented unit vectors $v,\widehat v$ with $c:=\widehat v^Tv\ge0$,
\[
  \sin^2(2\Theta(\widehat v,v))
  =(2c)^2(1-c^2).
\]
Since
\[
  \norm{\widehat v-v}^2=2-2c,
\]
the source rewrites this as
\begin{equation}
  \sin^2(2\Theta(\widehat v,v))
  =\frac14\norm{\widehat v-v}^2
    \bigl(2-\norm{\widehat v-v}^2\bigr)
    \bigl(4-\norm{\widehat v-v}^2\bigr). \tag{YWS-3}
\end{equation}
A bound on $\sin2\theta$ alone does not control $\sin\theta$ globally because $\sin2\theta$ is small both near $0$ and near $\pi/2$. An additional acute-overlap hypothesis, such as $c\ge1/\sqrt2$, would be needed.

\begin{corollary}[Rank-one form]
Fix $j$ and assume
\[
  \Delta_j:=\min(\lambda_{j-1}-\lambda_j,
                  \lambda_j-\lambda_{j+1})>0.
\]
For normalized eigenvectors $v,\widehat v$ corresponding to the $j$th decreasing eigenvalues,
\[
  \sin\Theta(\widehat v,v)
  \le
  \frac{2\norm{E}_{\op}}{\Delta_j}.
\]
If $\widehat v^Tv\ge0$, then
\[
  \norm{\widehat v-v}
  \le
  \frac{2^{3/2}\norm{E}_{\op}}{\Delta_j}.
\]
\end{corollary}

\section{The source's audit of statistical applications}

This section of the paper does not introduce a new perturbation theorem. It makes the following mathematical diagnosis of existing proof patterns.

\begin{enumerate}[leftmargin=2em]
\item In covariance, correlation, sparse PCA, factor-model, false-discovery, and leave-one-observation analyses, authors often apply a mixed-gap Davis--Kahan theorem and then use Weyl's inequality to recover a population-gap denominator. The population-gap theorem packages those two steps into one deterministic result.
\item Spectral clustering papers apply variants to adjacency matrices or graph Laplacians. Statements that appear to assume only a population gap may still hide a requirement that the perturbed and population matrices have the same number of eigenvalues in the selected interval. This is effectively a mixed spectral-separation condition.
\item The new theorem avoids that hidden interval-counting requirement by selecting corresponding ordered eigenblocks and imposing only the two exterior population gaps.
\end{enumerate}

These are literature-comparison claims rather than independent formal theorem statements. They motivate the theorem's interface and explain why its denominator is useful in probabilistic applications.

\section{Extension to rectangular matrices and singular subspaces}

\sourceanchor{Theorem 3 and equations (10)--(12) in the Appendix proof.}

Let $A,\widehat A\in\R^{p\times q}$ and set
\[
  D:=\widehat A-A.
\]
Let the singular values be
\[
  \sigma_1\ge\cdots\ge\sigma_{\min(p,q)},
  \qquad
  \widehat\sigma_1\ge\cdots\ge\widehat\sigma_{\min(p,q)}.
\]
Fix
\[
  1\le r\le s\le\rank(A),
  \qquad
  d=s-r+1,
\]
and define the squared population singular gap
\[
  \Delta_{\mathrm{sv}}
  :=\min(\sigma_{r-1}^2-\sigma_r^2,
         \sigma_s^2-\sigma_{s+1}^2)>0,
\]
with source conventions
\[
  \sigma_0^2:=+\infty,
  \qquad
  \sigma_{\rank(A)+1}^2:=-\infty.
\]
Let $V,\widehat V\in\R^{q\times d}$ contain the corresponding right singular vectors:
\[
  Av_j=\sigma_ju_j,
  \qquad
  \widehat A\widehat v_j=\widehat\sigma_j\widehat u_j.
\]

\begin{theorem}[Yu--Wang--Samworth singular-subspace theorem]
\[
  \norm{\sin\Theta(\widehat V,V)}_{\F}
  \le
  \frac{
    2(2\sigma_1+\norm{D}_{\op})
    \min(\sqrt d\norm{D}_{\op},\norm{D}_{\F})
  }{\Delta_{\mathrm{sv}}}.
\]
Moreover, some $\widehat O\in O(d)$ satisfies
\[
  \norm{\widehat V\widehat O-V}_{\F}
  \le
  \frac{
    2^{3/2}(2\sigma_1+\norm{D}_{\op})
    \min(\sqrt d\norm{D}_{\op},\norm{D}_{\F})
  }{\Delta_{\mathrm{sv}}}.
\]
The identical statements hold for the corresponding left singular-vector blocks $U,\widehat U$.
\end{theorem}

The paper presents this as a population-gap counterpart of Wedin's generalized $\sin\theta$ theorem.

\section{Appendix lemma: Frobenius compression}

\begin{lemma}[Orthogonal compression]
Let $M\in\R^{m\times n}$. If $U\in\R^{m\times p}$ and $W\in\R^{n\times q}$ have orthonormal columns, then
\[
  \norm{U^TMW}_{\F}\le\norm{M}_{\F}.
\]
If instead $U$ and $W$ have orthonormal rows, then
\[
  \norm{U^TMW}_{\F}=\norm{M}_{\F}.
\]
\end{lemma}

For the inequality, extend $U,W$ to square orthogonal matrices $[U\ U_1]$ and $[W\ W_1]$. Orthogonal invariance of the Frobenius norm gives
\[
  \norm{M}_{\F}
  =\norm{
     \begin{pmatrix}U^T\\U_1^T\end{pmatrix}
     M
     \begin{pmatrix}W&W_1\end{pmatrix}
   }_{\F}
  \ge\norm{U^TMW}_{\F}.
\]
For orthonormal rows, $UU^T=I_m$ and $WW^T=I_n$, so
\[
  \norm{U^TMW}_{\F}^2
  =\tr(U^TMWW^TM^TU)
  =\tr(MM^TUU^T)
  =\tr(MM^T)
  =\norm{M}_{\F}^2.
\]

\section{Proof of the population-gap theorem}

Set
\[
  \Lambda:=\diag(\lambda_r,\dots,\lambda_s),
  \qquad
  \widehat\Lambda:=\diag(\widehat\lambda_r,\dots,\widehat\lambda_s).
\]
The perturbed eigen-equation gives the residual identity
\[
  0=\widehat\Sigma\widehat V-
       \widehat V\widehat\Lambda
   =\Sigma\widehat V-
       \widehat V\Lambda
     +E\widehat V-
       \widehat V(\widehat\Lambda-\Lambda).
\]
Hence
\[
  R:=\widehat V\Lambda-
      \Sigma\widehat V
  =E\widehat V-
    \widehat V(\widehat\Lambda-\Lambda).
\]

\subsection{Two upper bounds for the residual}

Using the compression lemma and Weyl's eigenvalue inequality,
\begin{align}
  \norm{R}_{\F}
  &\le\norm{E\widehat V}_{\F}
      +\norm{\widehat V(\widehat\Lambda-\Lambda)}_{\F}\notag\\
  &\le\sqrt d\norm{E}_{\op}
      +\norm{\widehat\Lambda-\Lambda}_{\F}\notag\\
  &\le2\sqrt d\norm{E}_{\op}. \tag{YWS-4}
\end{align}
Alternatively, using Frobenius compression and Hoffman--Wielandt,
\begin{align}
  \norm{R}_{\F}
  &\le\norm{E}_{\F}
      +\norm{\widehat\Lambda-\Lambda}_{\F}\notag\\
  &\le2\norm{E}_{\F}. \tag{YWS-5}
\end{align}
Therefore
\[
  \norm{R}_{\F}
  \le2\min(\sqrt d\norm{E}_{\op},\norm{E}_{\F}).
\]

\subsection{Population gap gives the lower bound}

Choose $V_1\in\R^{p\times(p-d)}$ such that
\[
  P:=[V\ V_1]\in O(p),
\]
and
\[
  P^T\Sigma P=
  \begin{pmatrix}
    \Lambda&0\\0&\Lambda_1
  \end{pmatrix},
\]
where $\Lambda_1$ contains all eigenvalues outside the selected block.

Decompose $R$ in the $V\oplus V_1$ splitting. Orthogonality and Frobenius contraction give
\begin{align}
  \norm{R}_{\F}
  &\ge
  \norm{V_1V_1^T\widehat V\Lambda-
        V_1\Lambda_1V_1^T\widehat V}_{\F}\notag\\
  &\ge
  \norm{V_1^T\widehat V\Lambda-
        \Lambda_1V_1^T\widehat V}_{\F}. \tag{YWS-6}
\end{align}
Vectorization uses
\[
  \operatorname{vec}(ABC)
  =(C^T\otimes A)\operatorname{vec}(B),
\]
so
\begin{align}
 &\norm{V_1^T\widehat V\Lambda-
        \Lambda_1V_1^T\widehat V}_{\F}\notag\\
 &\quad=
 \norm{(\Lambda\otimes I_{p-d}-I_d\otimes\Lambda_1)
       \operatorname{vec}(V_1^T\widehat V)}\notag\\
 &\quad\ge
 \Delta\norm{\operatorname{vec}(V_1^T\widehat V)}. \tag{YWS-7}
\end{align}
Every diagonal entry of the Kronecker difference has magnitude at least $\Delta$ because one eigenvalue lies inside the selected population block and the other lies outside it.

Finally,
\begin{align*}
  \norm{\operatorname{vec}(V_1^T\widehat V)}^2
  &=\tr(\widehat V^TV_1V_1^T\widehat V)\\
  &=\tr((I-VV^T)\widehat V\widehat V^T)\\
  &=d-\norm{\widehat V^TV}_{\F}^2\\
  &=\norm{\sin\Theta(\widehat V,V)}_{\F}^2.
\end{align*}
Combining (YWS-4)--(YWS-7) proves (YWS-1).

\subsection{Procrustes alignment}

Choose $O_1,O_2\in O(d)$ so that
\[
  O_1^T\widehat V^TVO_2
  =\diag(\cos\theta_1,\dots,\cos\theta_d),
\]
and set
\[
  \widehat O:=O_1O_2^T.
\]
Then
\begin{align}
  \norm{\widehat V\widehat O-V}_{\F}^2
  &=2d-2\sum_{j=1}^d\cos\theta_j\notag\\
  &\le2d-2\sum_{j=1}^d\cos^2\theta_j\notag\\
  &=2\norm{\sin\Theta(\widehat V,V)}_{\F}^2. \tag{YWS-8}
\end{align}
Together with (YWS-1), this gives (YWS-2).

\section{Proof of the singular-subspace theorem}

Apply the symmetric theorem to the Gram matrices
\[
  A^TA,
  \qquad
  \widehat A^T\widehat A.
\]
Their selected eigenvectors are precisely the right singular vectors, and the eigenvalues are the squared singular values. Thus
\begin{equation}
  \norm{\sin\Theta(\widehat V,V)}_{\F}
  \le
  \frac{2\min\bigl(
      \sqrt d\norm{\widehat A^T\widehat A-A^TA}_{\op},
      \norm{\widehat A^T\widehat A-A^TA}_{\F}
    \bigr)}{\Delta_{\mathrm{sv}}}. \tag{YWS-9}
\end{equation}

Since
\[
  \widehat A^T\widehat A-A^TA
  =D^T\widehat A+A^TD,
\]
submultiplicativity gives
\begin{align}
  \norm{\widehat A^T\widehat A-A^TA}_{\op}
  &\le(\norm{\widehat A}_{\op}+\norm{A}_{\op})
       \norm{D}_{\op}\notag\\
  &\le(2\sigma_1+\norm{D}_{\op})\norm{D}_{\op}. \tag{YWS-10}
\end{align}
Here
\[
  \norm{\widehat A}_{\op}
  \le\norm{A}_{\op}+\norm{D}_{\op}
  =\sigma_1+\norm{D}_{\op}.
\]

For the Frobenius estimate, vectorization and Kronecker operator norms yield
\begin{align}
  \norm{\widehat A^T\widehat A-A^TA}_{\F}
  &\le
  \norm{(\widehat A^T\otimes I_q)
        \operatorname{vec}(D^T)}
  +\norm{(I_p\otimes A^T)\operatorname{vec}(D)}\notag\\
  &\le
  (\norm{\widehat A}_{\op}+\norm{A}_{\op})
  \norm{D}_{\F}\notag\\
  &\le
  (2\sigma_1+\norm{D}_{\op})\norm{D}_{\F}. \tag{YWS-11}
\end{align}
Substituting (YWS-10)--(YWS-11) into (YWS-9) proves the right-subspace bound. The Procrustes conclusion follows from (YWS-8). Applying the same argument to $AA^T$ and $\widehat A\widehat A^T$ proves the left-subspace version.

\section{Formalization map}

\begin{formalizationmap}
Principal angles & Singular values of $\widehat V^TV$, Frobenius sine distance, cross-complement identity. & \code{PrincipalAngles.lean}, \code{AlignedBasis.lean}, and the canonical subspace API used by \code{DavisKahanTheory/Statistics.lean}. \\
Population residual lower bound & $\Delta^2$ times cross-subspace energy is bounded by squared residual-column energy. & \code{YuWangSamworth.lean}: \code{sq_gap_mul_sum_cross_le_sum_sq_norm_residualColumn} and \code{sq_gap_mul_sum_cross_le_of_population_gap}. Fully proved. \\
Operator-norm numerator & Residual bounded by $2\sqrt d\norm{E}_{\op}$. & \code{YuWangSamworth.lean}: \code{sq_gap_mul_sum_cross_le_of_population_gap_opNorm}. Fully proved. \\
Sine theorem & Population-only gap bound for Frobenius principal-angle distance. & \code{YuWangSamworth.lean}: \code{sqrt_sum_cross_le_of_population_gap}; \code{DavisKahanTheory/Statistics.lean}: \code{yuWangSamworth_sinTheta_le}, \code{yuWangSamworth_intervalBlock_le}. Fully proved. \\
Procrustes step & Existence of aligned orthonormal bases with squared distance at most twice squared sine distance. & \code{AlignedBasis.lean}: \code{sum_sq_norm_aligned_le}; \code{Statistics.lean}: \code{exists_aligned_orthonormalBasis}, \code{yuWangSamworth_alignedBasis_le}. Fully proved. \\
Rank-one corollary & Sign-aligned eigenvector difference with constant $2^{3/2}$. & \code{Statistics.lean}: \code{yuWangSamworth_eigenvector_le}. Fully proved. \\
Rectangular extension & Right and left singular subspace sine bounds. & \code{DavisKahanTheory/SingularSubspace.lean}: \code{rightSingularSubspace_sinTheta_le}, \code{leftSingularSubspace_sinTheta_le}, and \code{singularSubspace_dilation_sinTheta_le}. Fully proved. \\
Frobenius compression & Orthogonal compression cannot increase Frobenius norm. & Replaced by intrinsic orthonormal-family and unitarily invariant norm lemmas rather than copied as a matrix-only helper. \\
\end{formalizationmap}

\section{Source-faithfulness and transcription audit}

\begin{enumerate}[leftmargin=2em]
\item The paper appeared in 2015, while the supplied source is the 2014 preprint transcription. The theorem numbering used here follows the transcription's labeled structure.
\item The transcription contains a duplicate, deprecated Davis--Kahan theorem and unfinished scratch calculations after its end-of-document marker. They are not part of the paper and are excluded from the reconstruction.
\item The baseline mixed-gap theorem's printed endpoint conventions should be checked against the official article before being used as a canonical index formula. The population-gap theorem and its proof are internally consistent.
\item The source uses finite real matrices. The Lean development generalizes the central symmetric-operator statements to finite-dimensional real or complex inner-product spaces and expresses subspaces intrinsically.
\item The source's SVD theorem is obtained through squared Gram matrices. The repository additionally contains a self-adjoint dilation route, which is mathematically equivalent but should be identified as a modern alternative rather than attributed to the paper.
\end{enumerate}

\begin{thebibliography}{8}
\bibitem{YWS} Y. Yu, T. Wang, and R. J. Samworth, \emph{A Useful Variant of the Davis--Kahan Theorem for Statisticians}, Biometrika 102 (2015), 315--323.
\bibitem{DavisKahan1970} C. Davis and W. M. Kahan, \emph{The Rotation of Eigenvectors by a Perturbation. III}, SIAM J. Numer. Anal. 7 (1970), 1--46.
\bibitem{Wedin1972} P.-\AA. Wedin, \emph{Perturbation bounds in connection with singular value decomposition}, BIT 12 (1972), 99--111.
\bibitem{HoffmanWielandt1953} A. J. Hoffman and H. W. Wielandt, \emph{The variation of the spectrum of a normal matrix}, Duke Math. J. 20 (1953), 37--39.
\bibitem{StewartSun1990} G. W. Stewart and J.-G. Sun, \emph{Matrix Perturbation Theory}, Academic Press, 1990.
\end{thebibliography}

\end{document}
